\clearpage
\phantomsection
\chapter*{题目索引}
\addcontentsline{toc}{chapter}{题目索引}
\par\medskip\noindent\textbf{第1章}\par
\BookIndexEntry{os-ex1-1}{例 1-1　进程越多，处理器利用率一定越高吗 【2022】}
\BookIndexEntry{os-ex1-2}{例 1-2　执行系统调用时，哪些操作属于软件 【2022】}
\BookIndexEntry{os-ex1-3}{例 1-3　微内核的收益与代价 【2023】}
\BookIndexEntry{os-ex1-4}{例 1-4　每次开机需要重新建立什么 【2022】}
\BookIndexEntry{os-1-review}{回顾与自测}
\par\medskip\noindent\textbf{第2章}\par
\BookIndexEntry{os-ex2-1}{例 2-1　主动让出处理器后去哪里 【2023】}
\BookIndexEntry{os-ex2-2}{例 2-2　先排除尚未到达者，再比较 【2017】}
\BookIndexEntry{os-ex2-3}{例 2-3　被抢占者的剩余时间必须保留 【2023】}
\BookIndexEntry{os-ex2-4}{例 2-4　一次到达不必增加一次进程调度 【2022】}
\BookIndexEntry{os-ex2-5}{例 2-5　二级反馈中，第二级按剩余量选择 【2019】}
\BookIndexEntry{os-ex2-6}{例 2-6　同一组到达事件，SRTF 与 RR 如何分别推进 〖2.2-006〗}
\BookIndexEntry{os-2-review}{自测：用相同时间线检查不同问题}
\par\medskip\noindent\textbf{第3章}\par
\BookIndexEntry{ex-err11}{例 3-1　共享缓冲区中的两种约束 〖2.3-011〗}
\BookIndexEntry{ex-err18}{例 3-2　同步信号量一定从 0 开始吗 〖2.3-018〗}
\BookIndexEntry{ex-2022}{例 3-3　哪些边需要跨线程通知 【2022】}
\BookIndexEntry{ex-2024}{例 3-4　同一个缓冲区为什么有不同的信号量方案 【2024】}
\BookIndexEntry{ex-2009}{例 3-5　消费者只能取特定类别时怎样计数 【2009】}
\BookIndexEntry{ex-2014}{例 3-6　连续取 10 件改变的是哪一层互斥 【2014】}
\BookIndexEntry{ex-2018}{例 3-7　条件变量的 wait 做什么 【2018】}
\BookIndexEntry{sec-selfcheck}{自测：先判断，再看答案}
\par\medskip\noindent\textbf{第4章}\par
\BookIndexEntry{os-ex4-1}{例 4-1　避免与检测需要的信息不同 【2015】}
\BookIndexEntry{os-ex4-2}{例 4-2　能先完成两人，仍然可能不安全 【2011】}
\BookIndexEntry{os-ex4-3}{例 4-3　把每个候选可用向量独立试一遍 〖2.4-026〗}
\BookIndexEntry{os-4-review}{自测：分别判断现在和最坏未来}
\par\medskip\noindent\textbf{第5章}\par
\BookIndexEntry{os-ex5-1}{例 5-1　页号可以不同，共享的页框必须相同 【2023】}
\BookIndexEntry{os-ex5-2}{例 5-2　最佳适应每次都重新比较当前空闲区 【2010】}
\BookIndexEntry{os-ex5-3}{例 5-3　一张下级表能覆盖多少页 【2010】}
\BookIndexEntry{os-ex5-4}{例 5-4　10 位字段不能按十六进制整位截断 【2019】}
\BookIndexEntry{os-ex5-5}{例 5-5　缺页次数与置换次数相差初装入次数 【2019】}
\BookIndexEntry{os-ex5-6}{例 5-6　缺页后，旧页框字段不能直接沿用 【2021】}
\BookIndexEntry{os-ex5-7}{例 5-7　窗口内有重复，集合只记一次 【2016】}
\BookIndexEntry{os-5-review}{自测：先写地址字段和状态，再计算}
\par\medskip\noindent\textbf{第6章}\par
\BookIndexEntry{os-ex6-1}{例 6-1　最后关闭，释放的是哪一层 【2023】}
\BookIndexEntry{os-ex6-2}{例 6-2　权限可以同时成立，不能只编码一种 【2017】}
\BookIndexEntry{os-ex6-3}{例 6-3　混合索引中，一次访问要经过几层 【2015】}
\BookIndexEntry{os-ex6-4}{例 6-4　链接文件修改一条记录，要把读和写都计入 〖4.2-025〗}
\BookIndexEntry{os-ex6-5}{例 6-5　同一 inode 不要求同一读写位置 【2017】}
\BookIndexEntry{os-ex6-6}{例 6-6　先定位位图块，再定位块内字节 【2015】}
\BookIndexEntry{os-6-review}{自测：沿名称与数据两条线检查}
\par\medskip\noindent\textbf{第7章}\par
\BookIndexEntry{os-ex7-1}{例 7-1　哪里需要知道具体设备参数 【2013】}
\BookIndexEntry{os-ex7-2}{例 7-2　首块要启动，末块要排空 【2011】}
\BookIndexEntry{os-ex7-3}{例 7-3　柱面、盘面、扇区是混合进位 〖5.3-016〗}
\BookIndexEntry{os-ex7-4}{例 7-4　根目录属于文件系统，不属于扇区编码 【2017】}
\BookIndexEntry{os-ex7-5}{例 7-5　方向决定先服务哪一侧 【2015】}
\BookIndexEntry{os-7-review}{自测：先找可重叠阶段，再核单位}
\clearpage
\phantomsection
\chapter*{方法索引}
\addcontentsline{toc}{chapter}{方法索引}
\par\medskip\noindent\textbf{第1章}\par
\BookIndexEntry{os-1-1}{操作系统管理什么，向程序提供什么}
\BookIndexEntry{os-c01-h13}{并发、共享、虚拟与异步怎样联系}
\BookIndexEntry{os-1-2}{批处理、分时与实时解决不同的问题}
\BookIndexEntry{os-1-3}{用户态、内核态与特权指令}
\BookIndexEntry{os-1-4}{中断和异常怎样把控制权交给内核}
\BookIndexEntry{os-1-5}{系统调用与库函数：调用一段代码，还是请求内核服务}
\BookIndexEntry{os-1-6}{从源程序到可运行映像：链接、装入与地址绑定}
\BookIndexEntry{os-1-7}{进程地址空间中，代码、数据、堆和栈各放什么}
\BookIndexEntry{os-1-8}{内核结构：把哪些功能放进同一保护域}
\BookIndexEntry{os-1-9}{开机以后，内核怎样获得控制权}
\BookIndexEntry{os-1-10}{虚拟机：客户操作系统也可以成为被管理者}
\par\medskip\noindent\textbf{第2章}\par
\BookIndexEntry{os-2-1}{进程与线程：资源单位和执行单位}
\BookIndexEntry{os-2-2}{就绪、运行和阻塞：先判断缺少什么}
\BookIndexEntry{os-2-3}{用户级、内核级线程与映射模型}
\BookIndexEntry{os-2-4}{创建、撤销与组织：PCB 为什么是管理入口}
\BookIndexEntry{os-2-5}{共享内存、消息、管道和信号}
\BookIndexEntry{os-2-6}{调度层次与评价指标}
\BookIndexEntry{os-2-7}{什么时刻可以调度，抢占改变了哪一步}
\BookIndexEntry{os-2-8}{沿时间线执行调度算法}
\BookIndexEntry{os-2-9}{多处理机调度：负载与缓存都有代价}
\BookIndexEntry{os-2-10}{上下文切换与模式切换}
\par\medskip\noindent\textbf{第3章}\par
\BookIndexEntry{sec-concepts}{临界区约束什么，为什么还需要同步}
\BookIndexEntry{os-c03-h37}{怎样判断互斥方案是否合格}
\BookIndexEntry{sec-locks}{软件互斥与锁：检查和占用必须连起来}
\BookIndexEntry{os-c03-h64}{Peterson 算法怎样处理双方同时申请}
\BookIndexEntry{sec-semaphore}{信号量：先说明它代表什么，再计算初值}
\BookIndexEntry{sec-precedence}{从前驱图到线程代码}
\BookIndexEntry{sec-buffer}{缓冲区：容量许可与互斥范围分别建立}
\BookIndexEntry{sec-readers}{多个读者怎样作为一组申请资源}
\BookIndexEntry{sec-philosophers}{哲学家进餐：申请多个资源时怎样断开等待环}
\BookIndexEntry{sec-monitor}{管程与条件变量：等待时必须让出改变条件的机会}
\BookIndexEntry{os-c03-h483}{wait 负责等待，不替调用者判断条件}
\BookIndexEntry{os-c03-h503}{通知以后，何时重新检查条件}
\BookIndexEntry{sec-method}{写出并发代码后，怎样检验}
\par\medskip\noindent\textbf{第4章}\par
\BookIndexEntry{os-4-1}{死锁、饥饿和不安全状态}
\BookIndexEntry{os-4-2}{死锁预防：在规则中破坏一个条件}
\BookIndexEntry{os-4-3}{银行家算法：可用不等于可以批准}
\BookIndexEntry{os-4-4}{检测、解除和恢复}
\par\medskip\noindent\textbf{第5章}\par
\BookIndexEntry{os-5-1}{逻辑地址、物理地址和可寻址范围}
\BookIndexEntry{os-5-2}{共享、保护、分配和回收}
\BookIndexEntry{os-5-3}{连续分配与动态分区：空闲总量不等于可分配块}
\BookIndexEntry{os-5-4}{分页：页号查表，偏移不变}
\BookIndexEntry{os-5-5}{TLB 与多级页表：空间与时间分别计算}
\BookIndexEntry{os-c05-h105}{TLB 缓存的是转换，不是文件内容}
\BookIndexEntry{os-5-6}{分段和段页式：先检查段内地址}
\BookIndexEntry{os-5-7}{请求分页：合法但不驻留时会发生什么}
\BookIndexEntry{os-5-8}{页框分配与置换范围}
\BookIndexEntry{os-5-9}{页面置换：必须跟踪实际状态}
\BookIndexEntry{os-5-10}{内存映射文件：把文件内容放入虚拟地址视图}
\BookIndexEntry{os-5-11}{工作集、抖动与平均访问时间}
\BookIndexEntry{os-c05-h238}{数组访问怎样变成页引用串}
\par\medskip\noindent\textbf{第6章}\par
\BookIndexEntry{os-6-2}{建立、打开、读写、关闭与删除}
\BookIndexEntry{os-6-3}{文件保护：用户类别与权限组合}
\BookIndexEntry{os-6-4}{逻辑结构：字节流和记录怎样被使用}
\BookIndexEntry{os-6-6}{目录和路径解析：逐层找到名字}
\BookIndexEntry{os-6-7}{硬链接、软链接与什么时候回收文件}
\BookIndexEntry{os-6-8}{文件系统的磁盘结构与内存结构}
\BookIndexEntry{os-6-9}{空闲空间：位图、空闲链和成组链接}
\BookIndexEntry{os-6-10}{虚拟文件系统与挂载}
\par\medskip\noindent\textbf{第7章}\par
\BookIndexEntry{os-7-1}{设备、控制器、接口与端口}
\BookIndexEntry{os-7-2}{轮询、中断和 DMA：谁搬数据，谁等待}
\BookIndexEntry{os-7-3}{I/O 软件层次与一次请求的完整路径}
\BookIndexEntry{os-7-4}{阻塞、非阻塞与异步接口}
\BookIndexEntry{os-7-5}{缓冲：把输入、复制与计算放到同一时间线}
\BookIndexEntry{os-7-6}{设备分配、逻辑名称与驱动接口}
\BookIndexEntry{os-7-7}{SPOOLing：先进入磁盘井，再使用独占设备}
\BookIndexEntry{os-7-8}{磁盘结构、地址与格式化}
\BookIndexEntry{os-7-9}{磁盘调度：把运动顺序和运动距离分开}
\BookIndexEntry{os-7-10}{SSD：没有寻道，不等于没有管理开销}
